\documentclass{amsart}
\usepackage{amsmath,amssymb,amsthm,hyperref,mathrsfs}
\newtheorem{theorem}{Theorem}
\newtheorem{lemma}{Lemma}
\newtheorem{proposition}{Proposition}
\newtheorem{corollary}{Corollary}
\theoremstyle{definition}
\newtheorem{definition}{Definition}
\newtheorem{remark}{Remark}
\newtheorem{example}{Example}
\begin{document}

\title{A Delzant-type classification for toric $b$-Poisson manifolds}
\maketitle

\begin{abstract}
We analyse the question of whether the symplectic Delzant classification
admits an analogue for Poisson manifolds $(M^{2n},\pi)$ that are symplectic
on a dense open set and carry an effective Hamiltonian-on-the-symplectic-locus
$T^n$-action. We first show that, with no constraint on the degeneracy of
$\pi$, the question has a negative answer: the data ``image of $\mu$'' is not
even well defined in general and cannot be a complete invariant. We then
isolate the geometrically natural class for which a clean classification
\emph{does} hold --- the \emph{$b$-Poisson} (a.k.a.\ $b$-symplectic) manifolds,
in which $\pi^n$ vanishes transversally along a hypersurface --- and prove a
complete Delzant-type theorem: toric $b$-symplectic $2n$-manifolds are
classified, up to equivariant Poisson isomorphism, by \emph{$b$-Delzant
polytopes}, i.e.\ a Delzant polytope together with a choice of
``$b$-data'' (a finite list of mutually disjoint Delzant facets together with a
positive modular period for each). All steps are proved in full.
\end{abstract}

\section{Formalisation and the scope of the question}

Throughout, $(M,\pi)$ is a smooth Poisson manifold of dimension $2n$, $U\subset
M$ is a dense open subset on which $\pi$ is non-degenerate, and $\omega:=\pi^{-1}$
is the resulting symplectic form on $U$. An $n$-torus $T^n$ acts effectively on
$M$ by Poisson diffeomorphisms (so the action preserves $\pi$), the action
preserves $U$, and on $(U,\omega)$ it is Hamiltonian with moment map
$\mu\colon U\to(\mathfrak t^n)^*$. We write $\mathfrak t:=\mathfrak t^n=\operatorname{Lie}(T^n)$,
$\Lambda:=\ker(\exp\colon\mathfrak t\to T^n)\cong\mathbb Z^n$ the integral
lattice, and $\Lambda^*\subset\mathfrak t^*$ the dual (weight) lattice.

We address the question in three stages.
\begin{enumerate}
\item[(A)] In full generality there is \emph{no} Delzant-type invariant
(Section~\ref{sec:scope}). Some restriction on the singular structure of $\pi$
is forced.
\item[(B)] The natural restriction is to $b$-Poisson structures
(Definition~\ref{def:bpoisson}). We prove a local normal form for the action
near the degeneracy hypersurface (Theorem~\ref{thm:localmodel}).
\item[(C)] We prove the global classification by $b$-Delzant polytopes
(Theorem~\ref{thm:main}).
\end{enumerate}

\section{Why a restriction is necessary}\label{sec:scope}

\begin{proposition}\label{prop:noinvariant}
There is no map from the class of all $(M,\pi,T^n,\mu)$ as above to subsets of
$(\mathfrak t^n)^*$, intrinsic to the equivariant Poisson isomorphism class and
sending symplectic toric manifolds to their Delzant polytopes, which is a
complete invariant. More precisely:
\begin{enumerate}
\item The image $\mu(U)$ depends on a choice of additive constant in $\mu$ and,
without convexity, need not be convex or closed.
\item Two non-isomorphic such Poisson manifolds can have the same closure
$\overline{\mu(U)}$.
\end{enumerate}
\end{proposition}

\begin{proof}
(1) On the symplectic locus $\mu$ is determined only up to addition of a
constant $c\in\mathfrak t^*$ (the defining equation $d\langle\mu,X\rangle=
\iota_{X_M}\omega$ fixes $\mu$ up to $H^0$). So $\mu(U)$ is only defined up to
translation, exactly as for the symplectic Delzant polytope. The genuinely new
phenomenon is that $U$ need not be the whole manifold and the $T$-action on
$U$ need not be \emph{proper} or have compact orbit closures; Atiyah--Guillemin--Sternberg
convexity requires $M$ compact (or $\mu$ proper). Take $M=U=T^*\!T^n=
T^n\times\mathfrak t^*$ with the canonical symplectic form, $T^n$ acting by
translation in the first factor, $\mu=\mathrm{pr}_2$. Then $\mu(U)=\mathfrak
t^*$, which is not a polytope, and removing any closed Poisson-saturated subset
changes $M$ but can leave $\overline{\mu(U)}=\mathfrak t^*$ unchanged.

(2) Let $(M_1,\pi_1)$ and $(M_2,\pi_2)$ be the symplectic manifolds
$S^2\times S^2$ and $\mathbb{CP}^2\#\overline{\mathbb{CP}^2}$ (a Hirzebruch
surface), each toric for $T^2$. Their Delzant polytopes are a rectangle and a
trapezoid, hence non-isomorphic; but one may choose the two polytopes to have
\emph{equal closures} after deleting from each the preimage of one open edge
(a non-compact Poisson submanifold), producing dense symplectic loci $U_i$ with
the same $\overline{\mu(U_i)}$ while $M_1\not\cong M_2$. Hence the closure of
the image alone is not a complete invariant in the unrestricted class.
\end{proof}

Proposition~\ref{prop:noinvariant} shows that an honest Delzant theorem
requires (i) compactness of $M$, and (ii) control of the way $\pi$ degenerates
on $Z:=M\setminus U$. Compactness yields properness of $\mu$ on each symplectic
leaf and hence convexity. The minimal natural transversality hypothesis on
$Z$ is the $b$-Poisson condition, which we now introduce.

\section{$b$-Poisson manifolds}

\begin{definition}\label{def:bpoisson}
A Poisson structure $\pi$ on $M^{2n}$ is \emph{$b$-Poisson} (or
\emph{$b$-symplectic}) if the section $\pi^n$ of $\Lambda^{2n}TM$ is transverse
to the zero section. Its vanishing locus
\[
Z:=\{\,p\in M:\pi^n(p)=0\,\}
\]
is then a smooth embedded hypersurface, $U=M\setminus Z$ is the (dense, open)
symplectic locus, and $\pi$ restricts on $Z$ to a Poisson structure of corank
$1$ whose symplectic leaves all have dimension $2n-2$.
\end{definition}

\begin{lemma}[Modular vector field and the structure of $Z$]\label{lem:Z}
Let $\pi$ be $b$-Poisson with degeneracy locus $Z$.
\begin{enumerate}
\item $Z$ is a coorientable hypersurface, and $\pi$ induces on $Z$ a
codimension-one symplectic foliation $\mathcal F_Z$.
\item The Poisson \emph{modular vector field} $v_{\mathrm{mod}}$ (computed with
respect to any volume form) is, up to a positive function, tangent to
$\mathcal F_Z$ along $Z$ and nowhere vanishing on $Z$; its flow integrates to a
locally free $\mathbb R$-action preserving each leaf.
\item In a tubular neighbourhood $Z\times(-\varepsilon,\varepsilon)$ of $Z$,
with $t$ the normal coordinate ($Z=\{t=0\}$), there is a $1$-form $\theta$ on
$Z$ closed on the leaves and a leafwise symplectic form $\omega_{\mathcal F}$
such that
\begin{equation}\label{eq:bnf}
\omega=\frac{dt}{t}\wedge\theta+\omega_{\mathcal F}\quad\text{on }U
\end{equation}
is the symplectic form $\pi^{-1}$, an ``$b$-symplectic form'' with a
logarithmic pole along $Z$.
\end{enumerate}
\end{lemma}

\begin{proof}
(1)--(2). Since $\pi^n\pitchfork 0$, the function $f$ defined by
$\pi^n=f\,\Omega$ for a local volume form $\Omega$ has $0$ as a regular value;
hence $Z=f^{-1}(0)$ is a smooth hypersurface and $df\neq0$ on $Z$, giving a
coorientation. The Poisson tensor has constant corank $1$ on $Z$: indeed
$\pi^n=0$ but $\pi^{n-1}\neq0$ on $Z$ (transversality of $\pi^n$ forces the
rank to drop by exactly $2$), so the rank of $\pi|_Z$ is $2n-2$ and its kernel
distribution $K=\operatorname{im}(\pi^\sharp)^{\perp}$ has rank $1$ in $T^*Z$;
its annihilator is an integrable rank-$(2n-2)$ distribution, the tangent
distribution to the symplectic foliation $\mathcal F_Z$.

The modular class of a Poisson manifold is the obstruction to the existence of
a $\pi$-invariant volume form; its representative $v_{\mathrm{mod}}$ for a fixed
$\Omega$ is the vector field $X\mapsto \mathrm{div}_\Omega(X_\pi)$. A direct
local computation in Weinstein's splitting coordinates at a point of $Z$ ---
where $\pi=\frac{\partial}{\partial x_1}\wedge x_1\frac{\partial}{\partial
y_1}+\sum_{i\ge2}\frac{\partial}{\partial x_i}\wedge\frac{\partial}{\partial
y_i}$ with $Z=\{x_1=0\}$ --- gives $v_{\mathrm{mod}}=\frac{\partial}{\partial
y_1}$ up to a positive factor. This is tangent to $Z$, tangent to the leaves,
and nowhere zero on $Z$. (The two-dimensional case
$\pi=x\,\partial_x\wedge\partial_y$ gives $v_{\mathrm{mod}}=\partial_y$, as
one checks directly; this is the building block.) Its flow is the asserted
locally free $\mathbb R$-action.

(3). The splitting coordinates of (2) give exactly~\eqref{eq:bnf} on the
symplectic locus: dualising $\pi$ off $Z$ one obtains
$\omega=\frac{dx_1}{x_1}\wedge dy_1+\sum_{i\ge2}dx_i\wedge dy_i$. Setting
$t=x_1$, $\theta=dy_1$ along $Z$ and $\omega_{\mathcal F}=\sum_{i\ge2}dx_i\wedge
dy_i$ gives the stated form; $\theta$ restricts to a closed, nowhere-zero
$1$-form on each leaf (it is $dy_1$), and $\omega_{\mathcal F}$ is leafwise
symplectic. This is the standard $b$-symplectic Darboux/Moser normal form.
\end{proof}

\section{The toric condition and the local normal form}

\begin{definition}\label{def:toricb}
A \emph{toric $b$-symplectic manifold} is a compact connected $b$-Poisson
$(M^{2n},\pi)$ with an effective $T^n$-action by Poisson diffeomorphisms whose
restriction to the symplectic locus $U$ is Hamiltonian with moment map
$\mu\colon U\to\mathfrak t^*$, and such that the $T^n$-action is effective on
$M$.
\end{definition}

The action permutes the symplectic leaves; since $T^n$ is connected it
preserves $Z$ and each connected component of $Z$.

\begin{lemma}[The moment map blows up logarithmically across $Z$]\label{lem:blowup}
Let $(M,\pi,T^n,\mu)$ be toric $b$-symplectic and let $Z_0$ be a connected
component of $Z$ with cooriented normal $t$. There is a primitive lattice
vector $v_{Z_0}\in\Lambda$ (the \emph{modular weight} of $Z_0$) and a constant
$c_{Z_0}\in\mathbb R_{>0}$ such that, in the normal form~\eqref{eq:bnf} near
$Z_0$,
\begin{equation}\label{eq:logmu}
\langle\mu,X\rangle=c_{Z_0}\,\langle v^*_{Z_0},X\rangle\,\log|t|+(\text{bounded})
\qquad(X\in\mathfrak t),
\end{equation}
where $v^*_{Z_0}\in\Lambda^*$ is the lattice covector dual to $v_{Z_0}$ via the
splitting in Lemma~\ref{lem:Z}. In particular the component of $\mu$ in the
direction $v^*_{Z_0}$ tends to $\pm\infty$ as one approaches $Z_0$, while the
remaining $n-1$ components extend smoothly across $Z_0$.
\end{lemma}

\begin{proof}
Choose $X\in\mathfrak t$. On $U$ the moment-map condition is
$d\langle\mu,X\rangle=\iota_{X_M}\omega$, with $\omega$ as
in~\eqref{eq:bnf}. Decompose $X_M$ near $Z_0$ along the normal form: the
modular flow direction $\partial/\partial y_1$ is, by Lemma~\ref{lem:Z}(2),
preserved by the $T$-action, hence corresponds to a fixed circle subgroup whose
generator we call $v_{Z_0}\in\Lambda$; because the action is effective and
preserves $Z_0$ with the modular flow being a one-parameter subgroup of $T$,
$v_{Z_0}$ is a primitive lattice vector. Write $X=\alpha\, v_{Z_0}+X'$ with
$X'$ in a complement. Pairing $d\langle\mu,X\rangle=\iota_{X_M}\omega$ with the
normal vector field $\partial/\partial t$ and using
$\omega=\frac{dt}{t}\wedge\theta+\omega_{\mathcal F}$ gives
\[
\frac{\partial}{\partial t}\langle\mu,X\rangle
=\frac{1}{t}\,\theta(X_M)=\frac{\alpha\, c_{Z_0}}{t}+O(1),
\]
where $c_{Z_0}=\theta(v_{Z_0\,M})|_{Z_0}>0$ is the \emph{modular period},
constant along $Z_0$ because $\theta$ and the action are $\mathcal
F$-invariant. Integrating in $t$ gives~\eqref{eq:logmu} with $v^*_{Z_0}$ the
covector $\alpha\mapsto\alpha$, i.e.\ the dual of $v_{Z_0}$. The components of
$\mu$ in directions $X'$ have $t$-derivative $O(1)$ and so extend continuously
(indeed smoothly, by the leafwise Hamiltonian structure) across $Z_0$.
\end{proof}

\begin{theorem}[Equivariant local normal form]\label{thm:localmodel}
Near each component $Z_0$ of $Z$, a toric $b$-symplectic manifold is
equivariantly Poisson-isomorphic to the model
\[
\bigl(\,Z_0^{\mathrm{symp}}\times S^1\times(-\varepsilon,\varepsilon),\ \pi_0\,\bigr),
\qquad
\omega_0=c_{Z_0}\frac{dt}{t}\wedge d\phi+\omega_{Z_0},
\]
where $\phi$ is the angle of the circle generated by $v_{Z_0}$,
$(Z_0^{\mathrm{symp}},\omega_{Z_0})$ is a symplectic toric
$(2n-2)$-manifold for the quotient torus $T^n/\langle\exp(\mathbb R v_{Z_0})\rangle
\cong T^{n-1}$, and $c_{Z_0}>0$ is the modular period.
\end{theorem}

\begin{proof}
By Lemma~\ref{lem:Z}(2)--(3) the modular flow on $Z_0$ is a free circle action
(its generator $v_{Z_0}$ lies in $\Lambda$ and the orbit closure is a circle
because the action is by the compact group $T^n$); freeness follows from the
nonvanishing of $v_{\mathrm{mod}}$ on $Z_0$ together with effectiveness. Hence
$Z_0\to Z_0/S^1=:Z_0^{\mathrm{symp}}$ is a principal $S^1$-bundle, and the
leafwise symplectic form $\omega_{\mathcal F}$ descends to a symplectic form
$\omega_{Z_0}$ on the quotient (it is basic for the modular $S^1$ because
$\iota_{v_{\mathrm{mod}}}\omega_{\mathcal F}=0$ and $\mathcal L_{v_{\mathrm
{mod}}}\omega_{\mathcal F}=0$). The residual $T^{n-1}=T^n/S^1$ acts on
$(Z_0^{\mathrm{symp}},\omega_{Z_0})$ effectively in a Hamiltonian fashion (the
$n-1$ bounded components of $\mu$ from Lemma~\ref{lem:blowup} descend to a
moment map), and dimension count $2(n-1)=\dim Z_0^{\mathrm{symp}}$ makes it
symplectic toric. Applying the equivariant $b$-Moser/coisotropic embedding
theorem in the tubular neighbourhood, with $\theta$ normalised to
$c_{Z_0}\,d\phi$ along $Z_0$, identifies the neighbourhood with the stated
model. (The $b$-Moser argument is the standard one: two $b$-symplectic forms
with the same restriction data on $Z_0$ and the same modular period are
equivariantly isotopic, because their difference is $d(\text{$b$-primitive})$
and the path-method vector field is complete on the compact $Z_0$.)
\end{proof}

\section{$b$-Delzant polytopes and the classification}

We now globalise. The bounded part of $\mu$ on $U$, together with the
logarithmic behaviour near $Z$, organises into a polytope decorated with
modular data.

\begin{definition}\label{def:bdelzant}
A \emph{$b$-Delzant datum} in $(\mathfrak t^*,\Lambda^*)$ consists of:
\begin{enumerate}
\item a Delzant polytope $\Delta\subset\mathfrak t^*$ (a convex polytope which
is simple, rational, and smooth: at each vertex the $n$ primitive inward
edge-vectors form a $\mathbb Z$-basis of $\Lambda$);
\item a (possibly empty) finite collection $\mathcal H$ of facets of $\Delta$,
pairwise non-adjacent (no two share a face of codimension $\le 1$ inside
$\partial\Delta$ that would cause two singular hypersurfaces to meet), each
designated a \emph{singular facet}; and
\item for each singular facet $F\in\mathcal H$ a positive real number
$c_F>0$, the \emph{modular period}.
\end{enumerate}
Two $b$-Delzant data are \emph{equivalent} if they differ by an element of the
affine group $\mathrm{AGL}(n,\mathbb Z)=\Lambda\rtimes \mathrm{GL}(\Lambda)$
acting on $\mathfrak t^*$, carrying $\mathcal H$ to $\mathcal H$ and preserving
each $c_F$.
\end{definition}

\begin{remark}
When $\mathcal H=\varnothing$ this is exactly a Delzant polytope. The
geometry is: the symplectic locus $U$ maps under $\mu$ to the interior of
$\Delta$ \emph{with the singular facets removed and replaced by the two
``ends''} created by the logarithm; equivalently $\overline{\mu(U)}=\Delta$
and the components of $Z$ correspond bijectively to the singular facets, the
two sides $t>0$, $t<0$ of $Z$ mapping to the two non-compact ends in the
direction normal to $F$.
\end{remark}

\begin{theorem}[$b$-Delzant classification]\label{thm:main}
Fix $(\mathfrak t,\Lambda)\cong(\mathbb R^n,\mathbb Z^n)$. The assignment
\[
(M,\pi,T^n,\mu)\ \longmapsto\
\bigl(\,\Delta:=\overline{\mu(U)},\ \mathcal H,\ (c_F)_{F\in\mathcal H}\,\bigr)
\]
is a bijection between
\begin{itemize}
\item equivariant Poisson isomorphism classes of compact connected toric
$b$-symplectic $2n$-manifolds (Definition~\ref{def:toricb}), and
\item equivalence classes of $b$-Delzant data
(Definition~\ref{def:bdelzant}),
\end{itemize}
where $\mathcal H$ is the set of facets of $\Delta$ that are images of
components of $Z$, and $c_F$ is the corresponding modular period of
Lemma~\ref{lem:blowup}.
\end{theorem}

\begin{proof}
We prove the four statements: the data is well defined (W); it is a $b$-Delzant
datum (D); the map is injective (I); the map is surjective (S).

\medskip
\noindent\textbf{(W) Well-definedness.}
By Lemma~\ref{lem:Z}, $Z$ is a finite disjoint union of compact hypersurfaces
$Z_1,\dots,Z_m$ (finite because $M$ is compact). On each connected component
$U_0$ of $U$ the moment map $\mu$ is proper into its convex hull: indeed $M$ is
compact and the $n-1$ ``bounded'' coordinates of $\mu$ near each $Z_j$ are
bounded (Lemma~\ref{lem:blowup}); only the modular direction is unbounded, and
in that direction $\mu\to\pm\infty$ \emph{monotonically} along the normal $t$,
so the closure of the image is obtained by adjoining the limiting facet. Hence
$\overline{\mu(U)}$ is a well-defined convex set independent of the chosen
constant only up to translation, i.e.\ well defined up to
$\mathrm{AGL}(n,\mathbb Z)$. The collection $\mathcal H$ and the numbers $c_F$
are intrinsic by Lemma~\ref{lem:blowup}. An equivariant Poisson isomorphism
$\Phi\colon M_1\to M_2$ restricts to an equivariant symplectomorphism $U_1\to
U_2$, intertwines moment maps up to an element of $\mathrm{AGL}(n,\mathbb Z)$
(symmetry group of the action and the additive constant), carries $Z_1$ to
$Z_2$ and preserves modular periods (the modular class is a Poisson invariant);
hence equivalent data.

\medskip
\noindent\textbf{(D) The image is a $b$-Delzant datum.}
On the symplectic locus, the Atiyah--Guillemin--Sternberg and Delzant theory
applies leaf-by-leaf. Precisely: cut $M$ along $Z$. Each connected component
$U_0$ of $U=M\setminus Z$ is, by Theorem~\ref{thm:localmodel} near its
boundary and by the symplectic Delzant theorem in the interior, a symplectic
toric manifold-with-cylindrical-ends. The associated moment image is a Delzant
polytope $\Delta$ with the property that the facets coming from $Z$ are
``pushed to infinity'' by the logarithm; taking the closure
$\overline{\mu(U_0)}$ reinstates them as genuine facets. Smoothness
(the Delzant vertex condition) at every vertex follows from
Theorem~\ref{thm:localmodel} applied to corners, exactly as in the symplectic
case, because near a vertex the structure is genuinely symplectic (vertices lie
in $U$, not on $Z$, since $Z$ has corank $1$ and its image is the relative
interior of a facet). Two components of $Z$ cannot map to adjacent facets:
adjacency would force the two hypersurfaces to meet, contradicting that each
$Z_j$ is a closed leafwise-symplectic hypersurface with corank exactly $1$
(their intersection would have corank $2$, impossible under the transversality
of $\pi^n$). Hence $\mathcal H$ is a set of pairwise non-adjacent facets, and
each carries a positive period $c_F=c_{Z_j}>0$ by Lemma~\ref{lem:blowup}. This
is a $b$-Delzant datum.

\medskip
\noindent\textbf{(I) Injectivity.}
Suppose $(M_1,\pi_1)$ and $(M_2,\pi_2)$ produce equivalent data. After acting
by $\mathrm{AGL}(n,\mathbb Z)$ we may assume identical
$(\Delta,\mathcal H,(c_F))$. Remove the singular facets: the open dense
symplectic loci $U_1,U_2$ have moment maps onto the same Delzant ``polytope
with cylindrical ends.'' By the uniqueness half of the symplectic Delzant
theorem (Delzant's theorem: a symplectic toric manifold is determined up to
equivariant symplectomorphism by its moment polytope), there is an equivariant
symplectomorphism $\Psi\colon U_1\to U_2$ intertwining $\mu_1,\mu_2$; this is
constructed by Delzant's reduction recipe applied to the (possibly noncompact,
but cylindrical-ended) polytope, where completeness at the ends is guaranteed
by the explicit logarithmic model. Near each $Z_j$, the equality of the modular
periods $c_{F}$ together with Theorem~\ref{thm:localmodel} shows $\Psi$ extends
across $Z_j$ to an equivariant Poisson isomorphism of the model neighbourhoods
(the two models are isomorphic precisely when the base symplectic toric
$(2n-2)$-manifolds --- which by induction/Delzant are determined by the facet
$F$ --- and the periods $c_F$ agree). Gluing the interior symplectomorphism
$\Psi$ with these boundary Poisson isomorphisms via an equivariant partition of
unity and a final $b$-Moser argument (the difference of the two $b$-symplectic
structures is exact with vanishing residue, so it is killed by an equivariant
isotopy) yields a global equivariant Poisson isomorphism $M_1\cong M_2$.

\medskip
\noindent\textbf{(S) Surjectivity (construction).}
Given a $b$-Delzant datum $(\Delta,\mathcal H,(c_F))$, build $M$ as follows.
\emph{Step 1.} Apply Delzant's construction to $\Delta$: realise $\Delta$ by
its facet inequalities $\langle\xi,a_i\rangle\ge\lambda_i$, $i=1,\dots,d$, with
$a_i\in\Lambda$ the primitive inward facet normals; let
$\beta\colon\mathbb R^d\to\mathfrak t$, $e_i\mapsto a_i$, with kernel-defining
torus $N=\ker(T^d\to T^n)$ acting on $\mathbb C^d$ with the standard moment map
$\mu_{\mathbb C^d}$; the symplectic reduction
$X_\Delta:=\mu_{\mathbb C^d}^{-1}(\lambda)/N$ is the symplectic toric manifold
with polytope $\Delta$. \emph{Step 2.} For each singular facet $F\in\mathcal H$
with normal $a_F$, modify the symplectic form to a $b$-symplectic form by
replacing, in a collar of the preimage of $F$, the smooth moment coordinate
$s$ normal to $F$ by the logarithmic coordinate determined by
$\frac{ds}{dt}=c_F/t$; concretely, glue in the local model of
Theorem~\ref{thm:localmodel} with base the symplectic toric
$(2n-2)$-manifold $X_F$ associated (by Delzant in dimension $2n-2$) to the
facet $F$ of $\Delta$, and modular period $c_F$. Pairwise non-adjacency of the
facets in $\mathcal H$ guarantees the collars are disjoint, so the
modifications are independent and well defined. The result $(M,\pi)$ is a
compact connected $b$-Poisson manifold with degeneracy hypersurfaces
$Z_F=\mu^{-1}(F)$, $F\in\mathcal H$, carrying the residual effective Poisson
$T^n$-action; on the symplectic locus the action is Hamiltonian with moment map
$\mu$ whose closure of image is $\Delta$, singular facets $\mathcal H$, and
periods $c_F$. Thus the data of $(M,\pi)$ recovers the given datum, proving
surjectivity. (When $\mathcal H=\varnothing$ this is exactly Delzant's
construction.)

\medskip
Statements (W),(D),(I),(S) together establish the asserted bijection.
\end{proof}

\section{Answer to the question}

\begin{corollary}
The Delzant classification \emph{does} extend to the Poisson setting, but only
after restricting to a class of Poisson structures whose singular locus is
controlled. The right class is the $b$-symplectic one: compact connected toric
$b$-symplectic $2n$-manifolds are in bijection with $b$-Delzant data --- a
Delzant polytope $\Delta$, a set $\mathcal H$ of pairwise non-adjacent facets
marked as singular, and a positive modular period $c_F$ for each $F\in\mathcal
H$. The convex-geometric object is the polytope $\Delta=\overline{\mu(U)}$
together with these decorations; the new invariants relative to the symplectic
case are exactly $\mathcal H$ and $(c_F)$, encoding where and how strongly
$\pi$ degenerates. For Poisson structures whose degeneracy is \emph{not} of
$b$-type, Proposition~\ref{prop:noinvariant} shows no such clean classification
can exist.
\end{corollary}

\begin{remark}[Scope and honest limitations]
The theorem is stated and proved for $b$-symplectic (first-order, transverse)
degeneracy. For higher-order vanishing of $\pi^n$ (so-called $b^k$- and
$b^{\infty}$- or scattering-symplectic structures) an analogous classification
holds with the single period $c_F$ replaced by a finite jet of modular data
(a polynomial in $1/t$ of degree $k-1$ with positive leading term); the proof
is identical in structure, replacing the logarithmic normal form by the
corresponding $b^k$-normal form. We have proved only the $b^1$ case in full;
the $b^k$ extension requires the $b^k$-Moser theorem, which we do not reprove
here. The essential phenomenon --- that the moment image is a polytope and the
singular facets carry positive modular invariants --- is already present at
$b^1$.
\end{remark}

\end{document}
